\newpage
\section{Introdução}
\label{sec:introducao}

Desde a pandemia de COVID-19, a biossegurança em ambientes fechados passou a receber maior atenção. Aspectos como a qualidade do ar, a densidade de ocupação e a limpeza das instalações são essenciais para a preservação da saúde dos indivíduos que frequentam esses espaços. Entretanto, a preocupação com a qualidade do ar ainda se concentra, em grande parte, no ambiente externo, desconsiderando o fato de que as pessoas passam aproximadamente 90\% do seu tempo em ambientes internos \cite{WHO2014MultipleExposure}. Nesse contexto, a qualidade do ar interior (do inglês Indoor Air Quality – IAQ) assume papel central no bem-estar humano e pode representar um risco significativo à saúde, uma vez que poluentes, vírus, bactérias, fungos e compostos orgânicos voláteis (VOCs) tendem a se acumular com maior intensidade em espaços fechados.

Sintomas como dores de cabeça, fadiga, cansaço excessivo e irritação nos olhos e pele estão entre os malefícios de um ambiente com baixa IAQ \cite{SAFE-IOT}. A qualidade do ar em um ambiente interno pode ser agravada devido a diversos fatores, como má circulação de ar, emissores de compostos voláteis indesejados/despercebidos e falta de manutenção do espaço. Além disso, como destacado anteriormente, indivíduos contemporâneos tendem a passar grande parte do seu dia em espaços internos. Dessa forma torna-se essencial monitorar a qualidade do ar nestes ambientes para garantir o bem-estar de seus ocupantes. 

Nesse contexto, aplicações de software contemporâneas no paradigma da Internet das Coisas (\textit{Internet of Things} - IoT) e Inteligência Artificial (IA), têm sido grandes aliados para monitorar ambientes internos e identificar situações críticas de qualidade de ar \cite{kelly2020internet}. Embora essas aplicações tenham surgido inicialmente em reação à pandemia de COVID-19, seu potencial vai além, pois podem ser adaptadas para diferentes cenários de emergências e desastres, oferecendo uma abordagem preventiva contra novas epidemias e pandemias \cite{MoniPaEp, SAFE-IOT, Oliveira2025HospitaisSeguros}.

Por isso, foi desenvolvido o sistema SAFE/UFRJ, que permite monitorar as condições de biossegurança \cite{UFRJBiosseguranca2020} em instalações (laboratórios, salas de aula e ambientes de convivência) por meio de dispositivos IoT \cite{rios2022exploring}. Estes dispositivos permitem a coleta de dados para indicar o número de pessoas em determinado ambiente, nível de IAQ, umidade e temperatura da instalação, apresentando essas informações em um \textit{dashboard}. O sistema ainda emite alertas quando níveis de risco (estabelecidos para cada instalação) são atingidos e possibilita o gerenciamento das instalações cadastradas, incluindo a solicitação de serviços de limpeza ou manutenção conforme o necessário \cite{SAFE-IOT}.

O sistema foi avaliado e evoluído com base nos requisitos elicitados em reuniões com os grupos de pesquisadores que compõem o projeto, sendo eles: desenvolvedores, analistas de requisitos e gerência. Desta forma, o sistema em seu estágio de desenvolvimento atual evidencia seu potencial no monitoramento de ambientes internos, tendo em vista que consegue coletar e apresentar informações acerca dos bioindicadores de espaços que possuem um dispositivo BIoT instalado, tudo isso, com uma periodicidade útil para auxiliar na tomada de decisões acerca destes espaços.

No entanto, existem funcionalidades que precisam ser revistas, avaliadas e evoluídas, para que desempenhem suas funções assim como especificado na documentação do sistema. Deste modo, o presente trabalho pretende descrever todo o processo realizado durante a evolução do módulo de contagem de pessoas presente no sistema SAFE/UFRJ. Iniciando-se na apresentação e explicação do módulo de contagem existente, expondo uma análise exploratória acerca dos dados coletados pelo sistema, para compreender o comportamento dos eventos observados e finalmente apresentando uma nova proposta com soluções para os problemas destacados e consecutivamente sua implementação no sistema já existente.

